%=============================================================================
% WAVE 4, ITERATION 19: P4 CONFIRMATION (DEAD) + FULL MATTER POWER SPECTRUM
%
% Agent: DFD W4i19 Agent 4 (P4 Power Transfer --- DEAD; repurposed)
% Model: Opus 4.6
% Date: 20 March 2026
%
% INHERITED STATE (i18):
%   - P4: DEAD. Three independent no-go arguments (W4i14). Final (W4i15).
%   - c_s = c CONFIRMED from corrected action. Full perturbation equation:
%     nabla^2 dpsi - (1/c^2) ddot{dpsi} = -(8piG/c^2) delta_rho
%   - K''(0) = 1 (Appendix Q subtraction). Background Delta = 0.
%   - G_eff(k) = G_N (k^2 + k_mu^2)/k^2, k_mu ~ 0.01 Mpc^{-1}
%   - psi does not cluster. Structure formation via baryons + G_eff only.
%
% MANDATE:
%   1. Confirm P4 still dead (1 paragraph)
%   2. MAIN: Full matter power spectrum P(k)_DFD from perturbation equation
%   3. Compute G_eff(k) and the psi-field Jeans scale k_psi
%   4. Quasi-static approximation validity domain
%   5. Estimate P(k)_DFD / P(k)_LCDM as function of k
%=============================================================================

\documentclass[aps,prd,twocolumn,nofootinbib]{revtex4-2}
\usepackage{amsmath,amssymb,bm,booktabs,xcolor}
\usepackage{graphicx}
\newcommand{\avec}{\bm{a}}
\newcommand{\aM}{a_0}
\newcommand{\astar}{a_*}
\newcommand{\ee}[1]{\times 10^{#1}}
\newcommand{\psifield}{\psi}
\newcommand{\dpsi}{\delta\psi}
\newcommand{\LCDM}{$\Lambda$CDM}
\newcommand{\Geff}{G_{\rm eff}}
\newcommand{\kmu}{k_\mu}
\newcommand{\kpsi}{k_\psi}
\newcommand{\kJ}{k_{\rm J}}
\newcommand{\ksh}{k_{\rm sh}}
\newcommand{\keq}{k_{\rm eq}}
\newcommand{\kdec}{k_{\rm dec}}
\newcommand{\kSilk}{k_{\rm Silk}}

\begin{document}

\title{W4i19 Agent 4: P4 Confirmation (DEAD) and\\
Full Matter Power Spectrum Calculation for DFD\\[6pt]
{\normalsize $G_{\rm eff}(k)$, Quasi-Static Approximation, and $P(k)_{\rm DFD}/P(k)_{\Lambda{\rm CDM}}$}}
\date{20 March 2026}
\maketitle

%=============================================================================
\section{P4 Status: DEAD (Confirmed)}
\label{sec:P4-dead}
%=============================================================================

P4 (Power Transfer Corridors) remains DEAD. The three independent no-go
arguments from W4i14 --- (1)~$c_\psi = c$, not $10^{-5}$~m/s, so the
scalar field cannot form a slow guided beam; (2)~no transverse confinement
mechanism exists in the DFD action ($W$ convexity forbids solitons,
W4i15); (3)~geometric $1/r^2$ dilution kills all long-range transfer
without confinement --- are strengthened by i18's confirmation that
$c_s = c$ at all epochs. The perturbation equation
$\nabla^2\dpsi - (1/c^2)\ddot\dpsi = -(8\pi G/c^2)\delta\rho$ is a
massless relativistic wave equation; there is no dispersive regime,
no slow group velocity, and no frequency band where guided propagation
could occur. All LCOE figures and distance-independence claims remain
invalidated. P4 is dead with no conceivable rescue pathway. No further
analysis is warranted.

%=============================================================================
%=============================================================================
\section{Setup: The Two-Field Perturbation System}
\label{sec:setup}
%=============================================================================
%=============================================================================

\subsection{The DFD perturbation equations}

We work with comoving coordinates, comoving wavenumber $k$, scale
factor $a(t)$, Hubble parameter $H = \dot a/a$. The two dynamical
fields are:
\begin{itemize}
\item $\dpsi(t,\mathbf{k})$: the psi-field perturbation (dimensionless).
\item $\delta_b(t,\mathbf{k}) \equiv \delta\rho_b/\bar\rho_b$: the
  baryon density contrast.
\end{itemize}

From i18 Agent~4, the psi-field perturbation equation in FRW is:
\begin{equation}
\boxed{
\ddot\dpsi + 3H\dot\dpsi + \frac{c^2 k^2}{a^2}\dpsi
= 8\pi G\,\bar\rho_b\,\delta_b.
}
\label{eq:psi-FRW}
\end{equation}

The baryon continuity + Euler equations, with the gravitational
potential $\Phi = -c^2\psi/2$, give:
\begin{equation}
\boxed{
\ddot\delta_b + 2H\dot\delta_b
= \frac{c^2 k^2}{2a^2}\dpsi.
}
\label{eq:baryon-growth}
\end{equation}

Equation~\eqref{eq:baryon-growth} states that baryons respond to
the psi-field gradient: $\nabla\Phi = -(c^2/2)\nabla\psi$, so
in Fourier space $k^2\Phi_k = (c^2/2)k^2\dpsi_k$ and the standard
growth equation $\ddot\delta + 2H\dot\delta = -k^2\Phi/a^2$ gives
Eq.~\eqref{eq:baryon-growth}.

\paragraph{Dimensional verification.}
\begin{itemize}
\item LHS of~\eqref{eq:psi-FRW}: $[\ddot\dpsi] = \text{s}^{-2}$
  (dimensionless field). $[3H\dot\dpsi] = \text{s}^{-1}\cdot\text{s}^{-1}
  = \text{s}^{-2}$. $[c^2 k^2 \dpsi/a^2] = (\text{m/s})^2
  \cdot\text{m}^{-2} = \text{s}^{-2}$. $\checkmark$
\item RHS of~\eqref{eq:psi-FRW}: $[8\pi G\bar\rho_b\delta_b]
  = (\text{m}^3\text{kg}^{-1}\text{s}^{-2})\cdot(\text{kg\,m}^{-3})
  = \text{s}^{-2}$. $\checkmark$
\item LHS of~\eqref{eq:baryon-growth}: $[\ddot\delta_b]
  = \text{s}^{-2}$. $\checkmark$
\item RHS of~\eqref{eq:baryon-growth}: $[c^2 k^2 \dpsi/(2a^2)]
  = \text{s}^{-2}$. $\checkmark$
\end{itemize}

\subsection{Connection to the standard Poisson equation}

In the quasi-static limit ($\ddot\dpsi \approx 0$, $3H\dot\dpsi
\approx 0$), Eq.~\eqref{eq:psi-FRW} reduces to:
\begin{equation}
\frac{c^2 k^2}{a^2}\dpsi = 8\pi G\,\bar\rho_b\,\delta_b,
\label{eq:QS-poisson}
\end{equation}
which is the Fourier-space Poisson equation
$\nabla^2\psi = -(8\pi G/c^2)\rho$ (with the sign conventions
matching our definition $\Phi = -c^2\psi/2$).

Substituting~\eqref{eq:QS-poisson} into~\eqref{eq:baryon-growth}:
\begin{equation}
\ddot\delta_b + 2H\dot\delta_b = 4\pi G\,\bar\rho_b\,\delta_b.
\label{eq:standard-growth}
\end{equation}
This is the standard CDM growth equation with $\Geff = G_N$.
In the quasi-static limit, DFD reproduces EXACTLY the Newtonian
baryon growth equation.


%=============================================================================
\section{The Quasi-Static Approximation: Validity Domain}
\label{sec:QSA}
%=============================================================================

This is the KEY question for structure formation.

\subsection{When is the QSA valid?}

The QSA drops the time derivatives $\ddot\dpsi$ and $3H\dot\dpsi$
relative to the gradient term $c^2 k^2\dpsi/a^2$ in
Eq.~\eqref{eq:psi-FRW}. This requires:
\begin{equation}
\frac{c^2 k^2}{a^2} \gg H^2 \quad\Longleftrightarrow\quad
k \gg \frac{aH}{c} \equiv \ksh.
\label{eq:QSA-condition}
\end{equation}

The quantity $\ksh = aH/c$ is the \textbf{sound horizon scale}
(equivalently, the comoving Hubble radius in units where $c_s = c$).

\paragraph{Numerical evaluation.}
At $z = 0$: $H_0 = 67.4$~km/s/Mpc $= 2.18\ee{-18}$~s$^{-1}$,
$a = 1$, $c = 3\ee{8}$~m/s:
\begin{equation}
\ksh(z=0) = \frac{H_0}{c} = \frac{2.18\ee{-18}}{3\ee{8}}
= 7.3\ee{-27}\;\text{m}^{-1}
\approx 2.2\ee{-4}\;\text{Mpc}^{-1}.
\label{eq:ksh-z0}
\end{equation}

At $z = 1100$ (recombination): $H(z_{\rm dec}) \approx
H_0\sqrt{\Omega_m}(1+z_{\rm dec})^{3/2}
\approx 67.4 \times 0.55 \times 3.65\ee{4}$~km/s/Mpc
$\approx 1.35\ee{6}$~km/s/Mpc $= 4.4\ee{-14}$~s$^{-1}$.
With $a = 1/1101$:
\begin{equation}
\ksh(z_{\rm dec}) = \frac{a\,H(z)}{c}
= \frac{(1/1101)\times 4.4\ee{-14}}{3\ee{8}}
\approx 1.3\ee{-25}\;\text{m}^{-1}
\approx 0.004\;\text{Mpc}^{-1}.
\end{equation}

Wait --- let me be more careful. The comoving sound horizon wavenumber
is:
\begin{equation}
\ksh = \frac{aH}{c_s} = \frac{aH}{c},
\end{equation}
but we need the COMOVING wavenumber, which means we compare
$c^2 k^2/a^2$ with $H^2$, giving:
\begin{equation}
k_{\rm phys} \gg \frac{H}{c}
\quad\Longleftrightarrow\quad
k_{\rm com} \gg \frac{aH}{c}.
\end{equation}

At any epoch, $aH/c$ in comoving Mpc$^{-1}$ is:
\begin{equation}
\ksh(a) = \frac{aH(a)}{c}
= \frac{H_0}{c}\,a\,E(a),
\end{equation}
where $E(a) = H(a)/H_0$. For matter domination, $E = a^{-3/2}$,
so $\ksh = (H_0/c)\,a^{-1/2}$.

At $z = 0$: $\ksh = H_0/c = 2.2\ee{-4}$~Mpc$^{-1}$.

At $z = 1100$: $\ksh = (H_0/c)(1+z)^{1/2}
= 2.2\ee{-4} \times 33.2 = 7.3\ee{-3}$~Mpc$^{-1}$.

\paragraph{Result.}
\begin{equation}
\boxed{
\text{QSA valid for } k \gg \ksh(z) \approx
\begin{cases}
7\ee{-3}\;\text{Mpc}^{-1} & (z = 1100), \\
2\ee{-4}\;\text{Mpc}^{-1} & (z = 0).
\end{cases}
}
\label{eq:QSA-validity}
\end{equation}

All cosmologically interesting scales ($k > 0.01$~Mpc$^{-1}$) are
well within the QSA domain at all post-recombination epochs. The
psi-field tracks the matter distribution quasi-statically (Poisson
equation) on all sub-Hubble scales.

\subsection{Physical interpretation}

When $c_s = c$, the psi-field ``sound crossing time'' across a mode
of wavelength $\lambda = 2\pi a/k$ is:
\begin{equation}
t_{\rm cross} = \frac{\lambda}{c_s} = \frac{2\pi a}{ck}.
\end{equation}
The Hubble time is $t_H = 1/H$. The QSA is valid when
$t_{\rm cross} \ll t_H$, i.e., the field adjusts much faster than
the universe expands. For $c_s = c$, this is identical to the condition
$k \gg aH/c$: the mode is inside the Hubble horizon.

\textbf{This is the crucial insight}: because $c_s = c$, the psi-field
adjusts at the speed of light. It cannot ``fall behind'' the matter
distribution on any sub-Hubble scale. It is enslaved to the matter
perturbations via the Poisson equation. This means:

\begin{enumerate}
\item The psi-field CANNOT cluster independently (it has no free
  growing modes --- only the forced response to matter).
\item The psi-field DOES respond to matter --- it provides the
  gravitational potential in which baryons grow.
\item The effective gravitational constant for baryon growth is
  $\Geff = G_N$ in the strict QSA (Newtonian regime).
\end{enumerate}


%=============================================================================
\section{The $\psi$-Field Jeans Scale $\kpsi$}
\label{sec:Jeans}
%=============================================================================

\subsection{Definition and derivation}

The Jeans wavenumber $\kJ$ for the psi-field is defined as the scale
at which the gradient pressure term equals the gravitational source.
From Eq.~\eqref{eq:psi-FRW}, in the matter era ($H^2 = 8\pi G
\bar\rho/3$), and seeking the threshold for growth vs.\ oscillation:

The homogeneous (source-free) psi equation is:
\begin{equation}
\ddot\dpsi + 3H\dot\dpsi + \frac{c^2 k^2}{a^2}\dpsi = 0.
\end{equation}

This has NO unstable modes for any $k > 0$. The gradient term
$c^2 k^2/a^2 > 0$ always provides a restoring force. The Jeans
wavenumber for the psi-field is:
\begin{equation}
\boxed{\kpsi = \infty.}
\label{eq:kpsi-inf}
\end{equation}

There is no gravitational instability threshold for the psi-field.
The field oscillates and decays on ALL sub-Hubble scales. This is
the defining characteristic of a $c_s = c$ field: it is
\emph{radiation-like} in its perturbation dynamics.

\subsection{The scale $\kmu$ is NOT the psi Jeans scale}

The scale $\kmu \approx 0.01$~Mpc$^{-1}$ appearing in the i17
$\Geff$ formula has a different origin. It comes from the NONLINEAR
$\mu(x)$ crossover in the DFD field equation, not from the linear
perturbation theory.

In the full nonlinear field equation:
\begin{equation}
\nabla\cdot\left[\mu\!\left(\frac{|\nabla\psi|}{\astar}\right)
\nabla\psi\right] = -\frac{8\pi G}{c^2}(\rho - \bar\rho),
\end{equation}
linearizing around a background gradient $|\nabla\bar\psi|/\astar
= x_0$ gives an effective $\Geff$ that depends on scale. The
characteristic scale is:
\begin{equation}
\kmu = \sqrt{\frac{8\pi G\bar\rho}{c^2\,\mu(x_0)}}\;\frac{a}{1},
\end{equation}
which for the Newtonian regime ($\mu \approx 1$) at $z = 0$ gives:
\begin{equation}
\kmu = \frac{a\sqrt{8\pi G\bar\rho_0}}{c}
= \frac{a\sqrt{3\Omega_m H_0^2}}{c}
= \frac{aH_0\sqrt{3\Omega_m}}{c}.
\end{equation}

At $z = 0$: $\kmu = (H_0/c)\sqrt{3\times 0.315}
= 2.2\ee{-4}\times 0.97 \approx 2.1\ee{-4}$~Mpc$^{-1}$.

\textbf{Wait}: This is $\kmu \approx \ksh$, NOT $0.01$~Mpc$^{-1}$.
Let me re-derive more carefully.

\subsection{Re-derivation of $\kmu$ from the coupled system}

The $\Geff(k)$ from i17 was derived by solving the COUPLED system
of Eqs.~\eqref{eq:psi-FRW} and~\eqref{eq:baryon-growth} in the
quasi-static limit. From the QSA:
\begin{equation}
\dpsi = \frac{8\pi G\bar\rho_b\delta_b\,a^2}{c^2 k^2}.
\end{equation}

Substituting into~\eqref{eq:baryon-growth}:
\begin{equation}
\ddot\delta_b + 2H\dot\delta_b
= \frac{c^2 k^2}{2a^2}\cdot\frac{8\pi G\bar\rho_b\delta_b\,a^2}{c^2 k^2}
= 4\pi G\bar\rho_b\delta_b.
\label{eq:growth-QSA}
\end{equation}

This gives $\Geff = G_N$ exactly. There is NO $k$-dependent
enhancement in the strict QSA of the LINEAR perturbation equation.

\textbf{Critical realization}: The i17 formula
$\Geff(k) = G_N(k^2 + \kmu^2)/k^2$ came from the NONLINEAR
$\mu(x)$ crossover, not from the linear perturbation theory.
In the linearized theory (where $\mu' = 1$, $\mu'' = 0$ around the
background), the response is simply Poisson.

The $\kmu$ enhancement arises when one accounts for the nonlinear
$\mu$-function in the quasi-static limit. Starting from the nonlinear
field equation and expanding around a LOCAL background (e.g., near a
galaxy or cluster), one obtains a scale-dependent response. But for
the COSMOLOGICAL background with $\nabla\bar\psi = 0$ and perturbations
in the Newtonian regime ($\mu \approx 1$), the linearized response
gives $\Geff = G_N$.

\textbf{This is a more pessimistic result than i17.}

\subsection{When does $\Geff \neq G_N$?}

The $\mu$-dependent $\Geff$ emerges in the DEEP-FIELD (MOND) regime
where $|\nabla\psi|/\astar \lesssim 1$. This requires:
\begin{equation}
|\nabla\psi| \lesssim \astar = \frac{2\aM}{c^2}
= \frac{2\times 1.2\ee{-10}}{9\ee{16}}
= 2.67\ee{-27}\;\text{m}^{-1}.
\end{equation}

For a perturbation with $\dpsi \sim \delta\Phi/c^2$ and
$\delta\Phi \sim G\bar\rho\,\delta\,\lambda^2$:
\begin{equation}
|\nabla\dpsi| \sim \frac{k}{a}\dpsi
\sim \frac{k}{a}\cdot\frac{4\pi G\bar\rho_b\delta_b\,a^2}{c^2 k^2}
= \frac{4\pi G\bar\rho_b\delta_b\,a}{c^2 k}.
\end{equation}

The condition $|\nabla\dpsi|/\astar \lesssim 1$ gives:
\begin{equation}
k \gtrsim \frac{4\pi G\bar\rho_b\delta_b\,a}{c^2\astar}
= \frac{4\pi G\bar\rho_b\delta_b\,a\,c^2}{c^2\cdot 2\aM}
= \frac{2\pi G\bar\rho_b\delta_b\,a}{\aM}.
\end{equation}

At $z = 0$: $\bar\rho_b = 4.2\ee{-28}$~kg/m$^3$,
$G = 6.67\ee{-11}$, $\aM = 1.2\ee{-10}$, $a = 1$, $\delta_b = 1$:
\begin{equation}
k_{\rm MOND} = \frac{2\pi\times 6.67\ee{-11}\times 4.2\ee{-28}}{1.2\ee{-10}}
= \frac{8.8\ee{-37}}{1.2\ee{-10}}
= 7.3\ee{-27}\;\text{m}^{-1}
\approx 2\ee{-4}\;\text{Mpc}^{-1}.
\end{equation}

This is the comoving Hubble scale! The MOND regime for cosmological
perturbations is at $k \lesssim 2\ee{-4}$~Mpc$^{-1}$, i.e.,
\textbf{super-Hubble scales}. On all sub-Hubble scales where we
measure the power spectrum, the perturbations are in the Newtonian
regime and $\Geff = G_N$.

\paragraph{Result.}
\begin{equation}
\boxed{
\Geff(k) = G_N \quad \text{for all } k > \ksh
\approx 2\ee{-4}\;\text{Mpc}^{-1}.
}
\label{eq:Geff-result}
\end{equation}

The i17 formula $\Geff = G_N(k^2 + \kmu^2)/k^2$ with $\kmu
= 0.01$~Mpc$^{-1}$ requires re-examination. That formula was
derived by assuming a cosmological MOND background, which is
inconsistent with the linearized treatment around a Newtonian
background (where $|\nabla\bar\psi| = 0$ and perturbations are
small). The correct linearized $\Geff$ is simply $G_N$.


%=============================================================================
\section{Full Matter Power Spectrum: $P(k)_{\rm DFD}$}
\label{sec:Pk}
%=============================================================================

\subsection{The DFD growth equation (post-recombination)}

With $\Geff = G_N$ from~\eqref{eq:Geff-result}, the baryon growth
equation after recombination ($z < 1100$) is:
\begin{equation}
\ddot\delta_b + 2H\dot\delta_b = 4\pi G\bar\rho_b\delta_b.
\label{eq:baryon-growth-final}
\end{equation}

This is IDENTICAL to the standard CDM growth equation in a universe
with matter fraction $\Omega_b = 0.049$ (baryons only). There is no
dark matter clustering component.

\subsection{Initial conditions at recombination}

At $z_{\rm dec} = 1100$, the baryon perturbation amplitude is set by
the baryon-photon acoustic oscillations:
\begin{equation}
\delta_b(k, z_{\rm dec}) = \delta_b^{\rm ac}(k)\,e^{-k^2/k_{\rm Silk}^2},
\end{equation}
where $\delta_b^{\rm ac}$ contains the BAO oscillations (peaks at
$k_n = n\pi/r_s$ with $r_s \approx 147$~Mpc the sound horizon) and
$\kSilk \approx 0.1$~Mpc$^{-1}$ is the Silk damping scale.

The amplitude of $\delta_b^{\rm ac}$ at recombination is:
\begin{equation}
|\delta_b(k, z_{\rm dec})| \sim 3\ee{-5}\quad
\text{(from Sachs-Wolfe)}.
\end{equation}

\subsection{Comparison with \LCDM{} initial conditions}

In \LCDM:
\begin{itemize}
\item CDM perturbations grow from matter-radiation equality
  ($z_{\rm eq} \approx 3400$). At $z_{\rm dec}$, CDM perturbations
  have grown by a factor $a_{\rm dec}/a_{\rm eq}
  = (1+z_{\rm eq})/(1+z_{\rm dec}) \approx 3.1$ relative to their
  primordial amplitude.
\item Baryons are still oscillating at $z_{\rm dec}$ (coupled to
  photons). After decoupling, they fall into the CDM potential wells.
  The ``catch-up'' happens within a few e-folds because the CDM wells
  are already deep ($\delta_{\rm CDM} \sim 3\delta_{\rm prim}$).
\item At late times, baryons and CDM have the same growth factor $D(a)$.
  The effective initial amplitude for the TOTAL matter perturbation
  ($\Omega_{\rm CDM}\delta_{\rm CDM} + \Omega_b\delta_b$) is dominated
  by CDM: $\delta_{\rm tot} \approx \delta_{\rm CDM}$.
\end{itemize}

In DFD:
\begin{itemize}
\item There is no CDM. The psi-field is smooth (quasi-static response
  only).
\item Baryons are the ONLY clustering component. At $z_{\rm dec}$,
  their amplitude is $\sim 3\ee{-5}$ with BAO oscillations and
  Silk damping.
\item After decoupling, baryons grow under $\Geff = G_N$ (Poisson
  equation from the quasi-static psi-field).
\item There are no pre-existing potential wells. The psi-field adjusts
  instantaneously to the baryon distribution but does not amplify it
  beyond the Newtonian response.
\end{itemize}

\subsection{Growth factor calculation}

The growing mode solution of~\eqref{eq:baryon-growth-final} in a
matter-dominated universe ($\Omega_m = 1$) is:
\begin{equation}
\delta_b \propto a \propto (1+z)^{-1}.
\end{equation}

The growth from $z_{\rm dec} = 1100$ to $z = 0$:
\begin{equation}
\frac{D(0)}{D(z_{\rm dec})} = \frac{1+z_{\rm dec}}{1+0} = 1101.
\end{equation}

But this is only valid if $4\pi G\bar\rho_b$ is the source. In a
universe with $\Omega_b = 0.049$ and $\Omega_\Lambda = 0.685$
($\Omega_m = 0.315$ in \LCDM), the growth equation is:
\begin{equation}
\ddot\delta_b + 2H\dot\delta_b = 4\pi G\bar\rho_b\delta_b
= \frac{3}{2}H^2\frac{\Omega_b}{\Omega_m}\,\Omega_m(a)\,\delta_b.
\end{equation}

\textbf{Critical point}: The factor $\Omega_b/\Omega_m
= 0.049/0.315 = 0.156$ appears because only baryons contribute to
the source (the psi-field is smooth). In \LCDM, the source is
$4\pi G\bar\rho_{\rm tot} = (3/2)H^2\Omega_m(a)$, which is
$\Omega_m/\Omega_b = 6.4\times$ larger.

However, this requires careful treatment. The background expansion
$H(a)$ is governed by the TOTAL energy density, including the psi-field
(which acts as dark energy/dark matter for the background). In DFD,
the background $H(z)$ matches \LCDM{} by construction (the distance-bias
framework, W4i14).

So the growth equation in DFD is:
\begin{equation}
\ddot\delta_b + 2H\dot\delta_b = 4\pi G\bar\rho_b\delta_b
= \frac{3}{2}H^2\Omega_b(a)\,\delta_b,
\label{eq:growth-DFD-full}
\end{equation}
where $\Omega_b(a) = \Omega_{b,0}\,a^{-3}/E^2(a)$ and $E(a)$
is the \LCDM{} expansion history.

In \LCDM, baryons + CDM grow together after decoupling:
\begin{equation}
\ddot\delta_m + 2H\dot\delta_m = 4\pi G\bar\rho_m\delta_m
= \frac{3}{2}H^2\Omega_m(a)\,\delta_m.
\label{eq:growth-LCDM}
\end{equation}

The ratio of the driving terms is:
\begin{equation}
\frac{4\pi G\bar\rho_b}{4\pi G\bar\rho_m}
= \frac{\Omega_b}{\Omega_m} = \frac{0.049}{0.315} = 0.156.
\label{eq:fb-ratio}
\end{equation}

\subsection{Growth suppression factor}

The growth equation $\ddot\delta + 2H\dot\delta = (3/2)H^2\Omega_X(a)\delta$
in the matter era has the growing mode $\delta \propto a^p$ where $p$
satisfies:
\begin{equation}
p(p-1) + \frac{3}{2}p = \frac{3}{2}\Omega_X.
\label{eq:growth-exponent}
\end{equation}

For $\Omega_X = 1$ (CDM in EdS): $p^2 + p/2 = 3/2$, giving $p = 1$.

For $\Omega_X = f_b \equiv \Omega_b/\Omega_m = 0.156$ (DFD):
$p^2 + p/2 = (3/2)\times 0.156 = 0.234$, giving
\begin{equation}
p = \frac{-1/2 + \sqrt{1/4 + 4\times 0.234}}{2}
= \frac{-0.5 + \sqrt{1.186}}{2}
= \frac{-0.5 + 1.089}{2}
= 0.295.
\end{equation}

\begin{equation}
\boxed{p_{\rm DFD} = 0.295 \quad\text{vs}\quad p_{\rm CDM} = 1.}
\label{eq:growth-exponents}
\end{equation}

The growth from $z_{\rm dec}$ to $z = 0$ (in the matter era):
\begin{align}
D_{\rm DFD} &\propto (1+z_{\rm dec})^{-p_{\rm DFD}}
= 1101^{-0.295}, \\
D_{\rm CDM} &\propto (1+z_{\rm dec})^{-1} = 1/1101.
\end{align}

Wait --- the growth factor is $D \propto a^p = (1+z)^{-p}$. From
$z_{\rm dec}$ to $z = 0$:
\begin{align}
\frac{D_{\rm DFD}(0)}{D_{\rm DFD}(z_{\rm dec})}
&= \left(\frac{a_0}{a_{\rm dec}}\right)^{p_{\rm DFD}}
= (1+z_{\rm dec})^{p_{\rm DFD}}
= 1101^{0.295} \approx 8.0, \label{eq:D-DFD}\\
\frac{D_{\rm CDM}(0)}{D_{\rm CDM}(z_{\rm dec})}
&= (1+z_{\rm dec})^{1} = 1101.
\label{eq:D-CDM}
\end{align}

\subsection{Initial condition difference}

In addition to the slower growth rate, DFD baryons start from a
DIFFERENT amplitude at $z_{\rm dec}$. In \LCDM:
\begin{itemize}
\item CDM has been growing since $z_{\rm eq} \approx 3400$.
  At $z_{\rm dec}$: $\delta_{\rm CDM} \approx
  \delta_{\rm prim}\times(1+z_{\rm eq})/(1+z_{\rm dec})
  \approx 3.1\,\delta_{\rm prim}$.
\item After decoupling, baryons fall into CDM wells within a few
  e-folds. The effective starting amplitude for the combined system
  is $\delta_m(z_{\rm dec}) \approx f_{\rm CDM}\delta_{\rm CDM}
  + f_b\delta_b^{\rm ac} \approx 0.84 \times 3.1\,\delta_{\rm prim}
  \approx 2.6\,\delta_{\rm prim}$ (dominated by CDM).
\end{itemize}

In DFD:
\begin{itemize}
\item The only clustering component is baryons.
  At $z_{\rm dec}$: $\delta_b \approx \delta_{\rm prim}$
  (the acoustic oscillation average --- peaks at
  $\sim 3\delta_{\rm prim}$, troughs near zero, and Silk damping
  suppresses small-scale modes).
\end{itemize}

The initial amplitude ratio:
\begin{equation}
\frac{\delta_b^{\rm DFD}(z_{\rm dec})}{\delta_m^{\rm LCDM}(z_{\rm dec})}
\approx \frac{1}{2.6} \approx 0.38.
\label{eq:IC-ratio}
\end{equation}

(This is an overestimate for DFD because Silk damping further
suppresses small-scale baryon modes. It is an underestimate for
\LCDM{} because the CDM transfer function enhances power on small
scales relative to large.)


%=============================================================================
\section{The Ratio $P(k)_{\rm DFD}/P(k)_{\rm \Lambda CDM}$}
\label{sec:ratio}
%=============================================================================

\subsection{Scale-independent estimate}

Combining the growth factor suppression and initial condition
difference:
\begin{align}
\frac{P_{\rm DFD}(k)}{P_{\rm CDM}(k)}
&= \left(\frac{\delta_b^{\rm DFD}(z_{\rm dec})}
  {\delta_m^{\rm CDM}(z_{\rm dec})}\right)^2
\times \left(\frac{D_{\rm DFD}}{D_{\rm CDM}}\right)^2 \notag\\
&\approx (0.38)^2 \times \left(\frac{8.0}{1101}\right)^2 \notag\\
&= 0.145 \times 5.3\ee{-5} \notag\\
&= 7.7\ee{-6}.
\label{eq:Pk-ratio-naive}
\end{align}

\begin{equation}
\boxed{
\frac{P_{\rm DFD}(k)}{P_{\rm CDM}(k)} \sim 10^{-5}
\quad\text{(scale-independent, naive)}.
}
\label{eq:Pk-ratio-boxed}
\end{equation}

This is a factor of $\sim 10^5$ deficit in the matter power spectrum
--- catastrophically low.

\subsection{Scale-dependent refinement}

The scale-independent estimate misses several effects:

\paragraph{(a) BAO oscillations in DFD initial conditions.}
The DFD baryon transfer function $T_b(k)$ at $z_{\rm dec}$ has
acoustic oscillations that \LCDM{}'s CDM transfer function $T_{\rm CDM}(k)$
does not. The CDM transfer function is a smooth turnover:
\begin{equation}
T_{\rm CDM}(k) \approx \frac{\ln(1 + 2.34\,q)}{2.34\,q}
\left[1 + 3.89\,q + (16.1\,q)^2 + (5.46\,q)^3
+ (6.71\,q)^4\right]^{-1/4},
\end{equation}
where $q = k/(h\,\Gamma)$ with $\Gamma = \Omega_m h\,
e^{-\Omega_b(1+\sqrt{2h/\Omega_m})}$ (Eisenstein-Hu without wiggles).

The DFD transfer function at $z_{\rm dec}$ is the BARYON transfer
function with full acoustic oscillations and Silk damping:
\begin{equation}
T_b^{\rm DFD}(k, z_{\rm dec}) \propto
\cos(kr_s)\,e^{-k^2/k_{\rm Silk}^2},
\end{equation}
where $r_s \approx 147$~Mpc and $\kSilk \approx 0.1$~Mpc$^{-1}$.

\paragraph{(b) Silk damping suppression.}
For $k > \kSilk \approx 0.1$~Mpc$^{-1}$, the DFD baryon
perturbations are exponentially suppressed by Silk damping:
\begin{equation}
\frac{P_{\rm DFD}}{P_{\rm CDM}} \propto e^{-2k^2/k_{\rm Silk}^2}
\quad \text{for } k > \kSilk.
\end{equation}
This makes the deficit WORSE on small scales.

\paragraph{(c) No CDM head-start bonus.}
In \LCDM, the CDM transfer function imprints the $z_{\rm eq}$
horizon scale $\keq \approx 0.01$~Mpc$^{-1}$ as the turnover.
In DFD, the relevant turnover is at $\kSilk$, which is larger.
Modes with $\keq < k < \kSilk$ are LESS suppressed in DFD
(baryon oscillations have moderate amplitude) but still lack the
CDM growth head-start.

\subsection{Scale-by-scale ratio}

\begin{center}
\begin{tabular}{lccc}
\toprule
Scale & $k$ (Mpc$^{-1}$) & $P_{\rm DFD}/P_{\rm CDM}$ & Comment \\
\midrule
Hubble & $10^{-3}$ & $\sim 10^{-5}$ & Both primordial; DFD has \\
 & & & no CDM growth \\
BAO peak & $0.05$ & $\sim 10^{-5}$--$10^{-4}$ & BAO peak in DFD; \\
 & & & no CDM turnover \\
Cluster & $0.1$ & $\sim 10^{-5}$ & Silk damping begins \\
Galaxy & $1$ & $\sim 10^{-9}$ & Silk exponential kills DFD \\
Dwarf gal.\ & $10$ & $\sim 10^{-40}$ & Silk damping extreme \\
\bottomrule
\end{tabular}
\end{center}

\paragraph{Note on $\Lambda$CDM late-time effects.}
The \LCDM{} growth is also suppressed relative to EdS at late times
($z < 1$) by $\Lambda$-domination. This affects both DFD and \LCDM{}
similarly (same $H(z)$), so the RATIO is determined primarily by the
$z_{\rm eq}$-to-$z_{\rm dec}$ head start and the growth exponent
difference.


%=============================================================================
\section{The $\sigma_8$ Prediction}
\label{sec:sigma8}
%=============================================================================

The variance $\sigma_8$ is:
\begin{equation}
\sigma_8^2 = \frac{1}{2\pi^2}\int_0^\infty dk\,k^2\,P(k)\,
|W(kR_8)|^2,
\end{equation}
where $R_8 = 8\,h^{-1}$~Mpc $\approx 12$~Mpc and $W$ is the
top-hat window function. The integral is dominated by
$k \sim 1/R_8 \approx 0.08$~Mpc$^{-1}$.

From the ratio estimate:
\begin{equation}
\sigma_8^{\rm DFD} \approx \sigma_8^{\rm CDM}
\times \sqrt{P_{\rm DFD}/P_{\rm CDM}}
\approx 0.81 \times \sqrt{10^{-5}}
\approx 0.81 \times 3.2\ee{-3}
= 2.6\ee{-3}.
\end{equation}

\begin{equation}
\boxed{\sigma_8^{\rm DFD} \approx 0.003
\quad\text{vs}\quad \sigma_8^{\rm obs} \approx 0.81.}
\end{equation}

This is a factor of $\sim 300\times$ too low (or $\sim 10^5$ in
power), consistent with the i18 Cosmo update's estimate of a
$10^{-6}$ deficit in power.

%=============================================================================
\section{Could Non-Linear Effects Help?}
\label{sec:nonlinear}
%=============================================================================

\subsection{Non-linear $\mu$ enhancement}

The MOND regime ($\mu < 1$) enhances gravity on scales where
$|\nabla\psi|/\astar < 1$. As shown in \S\ref{sec:Jeans}, this
only applies at $k \lesssim 2\ee{-4}$~Mpc$^{-1}$ for cosmological
perturbations. On galactic scales ($k \sim 1$--$10$~Mpc$^{-1}$),
the perturbation gradients are in the Newtonian regime.

However, at late times ($z < 2$), as structures go nonlinear,
individual halos enter the MOND regime in their outskirts. This
provides MOND-enhanced gravity within individual galaxy halos
(explaining flat rotation curves) but does NOT help the linear power
spectrum.

\subsection{External field effect and environment}

The EFE from large-scale structure gradients could modify the local
$\mu$ value on intermediate scales. However, this is a nonlinear
effect that cannot be captured by the linear perturbation theory.
The mochi\_class Boltzmann code would need to implement this via an
effective $\mu(k, z)$ function.

\subsection{Baryon feedback}

Post-recombination baryonic physics (star formation, AGN feedback,
reionization heating) affects the power spectrum on scales
$k > 1$~Mpc$^{-1}$ by 10--30\% in \LCDM. In DFD, these effects
would be MUCH more significant because baryons are the ONLY
clustering component, but they cannot provide the factor of $10^5$
needed to match observations.


%=============================================================================
\section{Comparison with AeST}
\label{sec:AeST}
%=============================================================================

The AeST theory (Skordis-Zlosnik 2021) solves this problem by
introducing a unit timelike vector field $A^\mu$ whose perturbations
behave as pressureless fluid (``ghost condensate''). In DFD, no such
field exists. The DFD psi-field is a single scalar with $c_s = c$;
it cannot split its clustering and MOND roles between different
degrees of freedom.

The structural lesson (confirmed again):
\begin{center}
\begin{tabular}{lcc}
\toprule
& MOND in galaxies & CDM in cosmology \\
\midrule
AeST & Scalar $\varphi$ (QS) & Vector $A^\mu$ (clusters) \\
DFD & Scalar $\psi$ (QS) & \textbf{Nothing} \\
\bottomrule
\end{tabular}
\end{center}


%=============================================================================
\section{Summary: The DFD Matter Power Spectrum}
\label{sec:summary}
%=============================================================================

\begin{enumerate}
\item The perturbation equation for $\dpsi$ is a massless wave equation
  with $c_s = c$. The psi-field Jeans wavenumber is
  $\kpsi = \infty$ --- no gravitational instability at any scale.

\item The quasi-static approximation is valid for ALL sub-Hubble modes
  ($k > aH/c$). The psi-field tracks the matter distribution via the
  Poisson equation.

\item In the QSA, the effective gravitational constant is
  $\Geff(k) = G_N$ for all cosmologically observable scales. The i17
  formula $\Geff = G_N(k^2 + \kmu^2)/k^2$ applies only in the MOND
  regime ($|\nabla\psi| \lesssim \astar$), which is at super-Hubble
  scales for cosmological perturbations. \textbf{This corrects the i17
  estimate.}

\item The baryon growth equation is $\ddot\delta_b + 2H\dot\delta_b
  = (3/2)H^2\Omega_b(a)\delta_b$, with growth exponent $p = 0.295$
  (vs.\ $p = 1$ for CDM). Combined with the missing CDM head-start
  ($z_{\rm eq}$ to $z_{\rm dec}$) and Silk damping:
  $P_{\rm DFD}/P_{\rm CDM} \sim 10^{-5}$ on scales
  $0.01 < k < 0.1$~Mpc$^{-1}$, worsening to $\sim 10^{-9}$ and below
  at $k > 1$~Mpc$^{-1}$.

\item $\sigma_8^{\rm DFD} \approx 0.003$, a factor $\sim 300\times$
  below the observed $\sigma_8 \approx 0.81$.
\end{enumerate}


%=============================================================================
\section*{TOP 5 FINDINGS}
\label{sec:findings}
%=============================================================================

\begin{table*}[t]
\centering
\caption{W4i19 Agent 4: Top 5 Findings, Ranked by Impact}
\begin{tabular}{@{}clp{11cm}@{}}
\toprule
\# & Category & Finding \\
\midrule
1 & \textbf{$P(k)_{\rm DFD}/P(k)_{\rm \Lambda CDM} \sim 10^{-5}$:
  CATASTROPHIC DEFICIT CONFIRMED AND QUANTIFIED} &
  The DFD matter power spectrum is suppressed by $\sim 10^5$ in power
  relative to observations on all cosmologically observable scales
  ($0.01 < k < 10$~Mpc$^{-1}$). This arises from THREE compounding
  effects: (a)~no CDM growth head-start ($z_{\rm eq}$ to $z_{\rm dec}$)
  reduces the initial amplitude by $\sim 3\times$;
  (b)~the growth exponent is $p = 0.295$ instead of $p = 1$
  (because the source is $4\pi G\bar\rho_b$ not $4\pi G\bar\rho_m$),
  giving a growth factor of 8 instead of 1101 from $z_{\rm dec}$;
  (c)~Silk damping exponentially suppresses baryon modes below
  $k_{\rm Silk}^{-1} \sim 10$~Mpc, making the deficit WORSE on
  small scales. The predicted $\sigma_8 \approx 0.003$ is $300\times$
  too low. This is \textbf{not rescuable} by any parameter adjustment
  within DFD's single-scalar framework. \\
\midrule
2 & \textbf{$\Geff(k) = G_N$ ON ALL OBSERVABLE SCALES: i17 $\kmu$
  FORMULA CORRECTED} &
  The i17 formula $\Geff = G_N(k^2 + \kmu^2)/k^2$ with $\kmu =
  0.01$~Mpc$^{-1}$ is incorrect for the linearized cosmological
  perturbation theory. In the linearized regime around the FRW
  background ($\nabla\bar\psi = 0$), the psi-field Poisson equation
  gives $\Geff = G_N$ with NO $k$-dependent enhancement. The MOND
  regime ($\mu < 1$, where $\Geff > G_N$) requires $|\nabla\psi|
  < \astar$, which corresponds to super-Hubble scales ($k < 2\ee{-4}$
  Mpc$^{-1}$) for cosmological perturbations. The $\kmu$ enhancement
  only applies to LOCAL nonlinear environments (galaxy/cluster halos),
  not to the linear power spectrum. This removes the partial rescue
  mechanism proposed in i17 and makes the deficit MORE severe. \\
\midrule
3 & \textbf{QSA VALID ON ALL SUB-HUBBLE SCALES: $\psi$ IS
  POISSON-SLAVED} &
  Because $c_s = c$, the quasi-static approximation holds for all
  modes with $k > aH/c \approx 7\ee{-3}$~Mpc$^{-1}$ at $z = 1100$
  and $k > 2\ee{-4}$~Mpc$^{-1}$ at $z = 0$. The psi-field responds
  instantaneously (at the speed of light) to the matter distribution.
  It is enslaved to baryons via the Poisson equation $\nabla^2\psi
  = -(8\pi G/c^2)\rho_b$, with no independent dynamics. This
  simultaneously (a)~explains why MOND works in galaxies (the QS
  psi-field provides the modified potential) and (b)~prevents the
  psi-field from acting as a CDM substitute (it has no free growing
  modes). \\
\midrule
4 & \textbf{$\kpsi = \infty$: THE PSI-FIELD HAS NO JEANS SCALE} &
  The psi-field perturbation equation is a massless wave equation with
  no gravitational instability at any finite wavelength. The Jeans
  wavenumber is formally infinite. This is the most fundamental
  reason why the psi-field cannot cluster: every sub-Hubble mode
  oscillates and decays rather than growing. The question ``what
  modes can $\psi$ enhance?'' has the answer: NONE, through
  self-clustering. The psi-field enhances gravity only through the
  quasi-static (Poisson) response, which gives $\Geff = G_N$ in
  the linear cosmological regime. \\
\midrule
5 & \textbf{P4 STILL DEAD (7th CONFIRMATION)} &
  The relativistic wave equation $c_s = c$ confirmed by i18 reinforces
  all three P4 no-go arguments. No dispersive regime, no slow group
  velocity, no guided propagation. 7th consecutive iteration confirming
  DEAD status. No further analysis warranted. \\
\bottomrule
\end{tabular}
\end{table*}


%=============================================================================
\section*{Implications for the DFD Programme}
%=============================================================================

The $10^{-5}$ power spectrum deficit is the most severe quantitative
test DFD has ever failed. The key distinction from previous estimates:

\begin{itemize}
\item \textbf{i17}: Estimated $P_{\rm DFD}/P_{\rm CDM} \sim 0.1$
  ($10\times$ deficit) using $\Geff(k) = G_N(k^2 + \kmu^2)/k^2$ as
  a rescue mechanism.

\item \textbf{i18}: Estimated $\sim 10^{-6}$ deficit but noted
  uncertainty in the $\Geff$ formula.

\item \textbf{i19 (this work)}: Confirms $\sim 10^{-5}$ deficit and
  shows that the $\Geff > G_N$ rescue mechanism does NOT operate on
  cosmologically observable scales (the MOND regime is super-Hubble).
  The growth exponent calculation ($p = 0.295$) is the key new result.
\end{itemize}

The resolution options remain as identified in i18:
\begin{enumerate}
\item \textbf{AeST-like vector extension}: Add a unit timelike vector
  $A^\mu$ whose perturbations cluster like CDM. Most promising but
  requires fundamental revision of DFD.
\item \textbf{Modified $K(\Delta)$}: Engineer $K$ so that the
  perturbation sound speed varies with epoch. Likely introduces
  instabilities.
\item \textbf{Accept DFD as galactic-scale theory}: Adopt the
  historical MOND position that dark matter is needed for cosmology.
\end{enumerate}

\textbf{mochi\_class remains existentially urgent}: The numerical
Boltzmann code would (a)~confirm or correct the $10^{-5}$ estimate,
(b)~test whether nonlinear $\mu$ effects provide any rescue at
intermediate scales, and (c)~compute the full CMB and BAO predictions
with proper initial conditions.

\end{document}
